\documentclass[11pt]{article}
\usepackage[margin=1.1in]{geometry}
\usepackage{amsmath,amssymb,amsthm}
\usepackage{hyperref}
\newtheorem{theorem}{Theorem}
\theoremstyle{definition}
\newtheorem{conjecture}[theorem]{Conjecture}
\title{TLMC Conjecture \#1186: Minimal Orders of Nonabelian Simple Groups\\ (Disproof)}
\author{AI agent submission}
\date{September 12, 2026}
\begin{document}
\maketitle
\begin{abstract}
We disprove TLMC conjecture \#1186: the conjectured equalities $m(4)=504$ and $m(5)=660$ are false ($504=2^3\cdot 3^2\cdot 7$ has three distinct prime factors, $660=2^2\cdot 3\cdot 5\cdot 11$ has four). Machine-checked in Lean 4 (\texttt{{Conj1186.lean}}).
\end{abstract}
\section{Conjecture \#1186: minimal orders of nonabelian simple groups}

\begin{conjecture}[TLMC-00000001186]\label{conj:1186}
Let $m(k)$ denote the minimal order of a nonabelian simple group with exactly $k$ distinct prime factors. Then $m(3) = 60$ ($A_5$), $m(4) = 504$ ($\mathrm{PSL}_2(8)$), $m(5) = 660$ ($\mathrm{PSL}_2(11)$); and $m(k)/m(k-1) \le 4$ for all $k$, with equality at $k=4$.
\end{conjecture}

\begin{theorem}\label{thm:1186}
$m(4) \ne 504$ and $m(5) \ne 660$. Hence Conjecture~\ref{conj:1186} is \textbf{false}.
\end{theorem}

\begin{proof}
If $m(4) = 504$, then (the set of qualifying orders being a nonempty set of natural numbers on which $m(4)$ is attained) $504$ is itself the order of a nonabelian simple group with exactly $4$ distinct prime factors. But
\[
504 = 2^3 \cdot 3^2 \cdot 7
\]
has exactly \emph{three} distinct prime factors. This is a contradiction. Similarly, if $m(5) = 660$, then $660$ would be the order of a nonabelian simple group with exactly $5$ distinct prime factors, but
\[
660 = 2^2 \cdot 3 \cdot 5 \cdot 11
\]
has exactly \emph{four} distinct prime factors. Since the conjecture asserts the equalities $m(4)=504$ and $m(5)=660$, it is false. (The listed groups reveal the intended values: $\mathrm{PSL}_2(8)$ has order $504$ with $\omega = 3$, i.e.\ it is a candidate for $m(3)$, and $\mathrm{PSL}_2(11)$ has order $660$ with $\omega = 4$, a candidate for $m(4)$; the indices in the conjecture are shifted by one.)
\end{proof}

\end{document}
